% GENERATED FILE. Edit canonical lessons, then run npm run generate:books.
% canonical-source-hash: fnv1a64:1492869a448b2877
% canonical-lessons: KA-C239-acarisu, KA-C239-alankarisu, KA-C239-hancu, KA-C239-bheti-kodu, KA-C239-appiko

\chapter{To celebrate, To decorate, To share out, To pay a visit, To hug}
\label{ch:ka-to-celebrate-to-decorate-to-share-out-to-pay-a-visit-to-hug}
\hlchaptermodality{239}

\begin{chapteropening}
\textbf{By the end of this chapter:} \emph{I can say to celebrate, to decorate, to share out, to pay a visit and to hug.}

Complete the last lesson of chapter 239: I can say to celebrate, to decorate, to share out, to pay a visit and to hug.
\end{chapteropening}

\section[\texorpdfstring{\kn{ಆಚರಿಸು} (ācarisu)}{ಆಚರಿಸು (ācarisu)}]{\kn{ಆಚರಿಸು} (ācarisu) — to celebrate}

\label{lesson:KA-C239-acarisu}



\begin{quote}
Before the new one: say the Kannada for to refuse, then the Kannada for to earn.
\end{quote}

\subsection*{You'll want to know: \kn{ಆಚರಿಸು}}
\textbf{\kn{ಆಚರಿಸು}} — \emph{ācarisu} — \textquotedblleft{}to celebrate\textquotedblright{}.

\begin{grammarlens}[title={What you've built}]
One more word for this chapter.
\end{grammarlens}

\subsection*{Your turn}
\begin{itemize}
\raggedright
  \item \emph{Say it:} \emph{ācarisu}
  \item \emph{Say it:} \emph{ācarisu}, once more
  \item \emph{Say it:} say \emph{ācarisu}
  \item \emph{Recall:} say the Kannada for to refuse, then the Kannada for to earn, then say \emph{ācarisu} again
\end{itemize}

\begin{tcolorbox}[breakable,colback=teal!4,colframe=teal!35!black,title={Before you move on}]
What does \kn{ಆಚರಿಸು} mean? (\textquotedblleft{}To celebrate\textquotedblright{}.) Say it once more.
\end{tcolorbox}
\section[\texorpdfstring{\kn{ಅಲಂಕರಿಸು} (alaṅkarisu)}{ಅಲಂಕರಿಸು (alaṅkarisu)}]{\kn{ಅಲಂಕರಿಸು} (alaṅkarisu) — to decorate}

\label{lesson:KA-C239-alankarisu}



\begin{quote}
Before the new one: say the Kannada for to earn, then the Kannada for to celebrate.
\end{quote}

\subsection*{You'll want to know: \kn{ಅಲಂಕರಿಸು}}
\textbf{\kn{ಅಲಂಕರಿಸು}} — \emph{alaṅkarisu} — \textquotedblleft{}to decorate\textquotedblright{}.

\begin{grammarlens}[title={What you've built}]
One more word for this chapter.
\end{grammarlens}

\subsection*{Your turn}
\begin{itemize}
\raggedright
  \item \emph{Say it:} \emph{alaṅkarisu}
  \item \emph{Say it:} \emph{alaṅkarisu}, once more
  \item \emph{Say it:} say \emph{alaṅkarisu}
  \item \emph{Recall:} say the Kannada for to earn, then the Kannada for to celebrate, then say \emph{alaṅkarisu} again
\end{itemize}

\begin{tcolorbox}[breakable,colback=teal!4,colframe=teal!35!black,title={Before you move on}]
What does \kn{ಅಲಂಕರಿಸು} mean? (\textquotedblleft{}To decorate\textquotedblright{}.) Say it once more.
\end{tcolorbox}
\section[\texorpdfstring{\kn{ಹಂಚು} (hañcu)}{ಹಂಚು (hañcu)}]{\kn{ಹಂಚು} (hañcu) — to share out, to distribute}

\label{lesson:KA-C239-hancu}



\begin{quote}
Before the new one: say the Kannada for to celebrate, then the Kannada for to decorate.
\end{quote}

\subsection*{You'll want to know: \kn{ಹಂಚು}}
\textbf{\kn{ಹಂಚು}} — \emph{hañcu} — \textquotedblleft{}to share out, to distribute\textquotedblright{}.

\begin{grammarlens}[title={What you've built}]
One more word for this chapter.
\end{grammarlens}

\subsection*{Your turn}
\begin{itemize}
\raggedright
  \item \emph{Say it:} \emph{hañcu}
  \item \emph{Say it:} \emph{hañcu}, once more
  \item \emph{Say it:} say \emph{hañcu}
  \item \emph{Recall:} say the Kannada for to celebrate, then the Kannada for to decorate, then say \emph{hañcu} again
\end{itemize}

\begin{tcolorbox}[breakable,colback=teal!4,colframe=teal!35!black,title={Before you move on}]
What does \kn{ಹಂಚು} mean? (\textquotedblleft{}To share out, to distribute\textquotedblright{}.) Say it once more.
\end{tcolorbox}
\section[\texorpdfstring{\kn{ಭೇಟಿ} \kn{ಕೊಡು} (bhēṭi koḍu)}{ಭೇಟಿ ಕೊಡು (bhēṭi koḍu)}]{\kn{ಭೇಟಿ} \kn{ಕೊಡು} (bhēṭi koḍu) — to pay a visit}

\label{lesson:KA-C239-bheti-kodu}



\begin{quote}
Before the new one: say the Kannada for to decorate, then the Kannada for to share out.
\end{quote}

\subsection*{You'll want to know: \kn{ಭೇಟಿ} \kn{ಕೊಡು}}
\textbf{\kn{ಭೇಟಿ} \kn{ಕೊಡು}} — \emph{bhēṭi koḍu} — \textquotedblleft{}to pay a visit\textquotedblright{}.

\begin{grammarlens}[title={What you've built}]
One more word for this chapter.
\end{grammarlens}

\subsection*{Your turn}
\begin{itemize}
\raggedright
  \item \emph{Say it:} \emph{bhēṭi koḍu}
  \item \emph{Say it:} \emph{bhēṭi koḍu}, once more
  \item \emph{Say it:} say \emph{bhēṭi koḍu}
  \item \emph{Recall:} say the Kannada for to decorate, then the Kannada for to share out, then say \emph{bhēṭi koḍu} again
\end{itemize}

\begin{tcolorbox}[breakable,colback=teal!4,colframe=teal!35!black,title={Before you move on}]
What does \kn{ಭೇಟಿ} \kn{ಕೊಡು} mean? (\textquotedblleft{}To pay a visit\textquotedblright{}.) Say it once more.
\end{tcolorbox}
\section[\texorpdfstring{\kn{ಅಪ್ಪಿಕೊ} (appiko)}{ಅಪ್ಪಿಕೊ (appiko)}]{\kn{ಅಪ್ಪಿಕೊ} (appiko) — to hug}

\label{lesson:KA-C239-appiko}



\begin{quote}
Before the new one: say the Kannada for to share out, then the Kannada for to pay a visit.
\end{quote}

\subsection*{You'll want to know: \kn{ಅಪ್ಪಿಕೊ}}
\textbf{\kn{ಅಪ್ಪಿಕೊ}} — \emph{appiko} — \textquotedblleft{}to hug\textquotedblright{}.

\begin{grammarlens}[title={What you've built}]
One more word for this chapter.
\end{grammarlens}

\subsection*{Your turn}
\begin{itemize}
\raggedright
  \item \emph{Say it:} \emph{appiko}
  \item \emph{Say it:} \emph{appiko}, once more
  \item \emph{Say it:} say \emph{appiko}
  \item \emph{Recall:} say the Kannada for to share out, then the Kannada for to pay a visit, then say \emph{appiko} again
\end{itemize}

\begin{tcolorbox}[breakable,colback=teal!4,colframe=teal!35!black,title={Before you move on}]
What does \kn{ಅಪ್ಪಿಕೊ} mean? (\textquotedblleft{}To hug\textquotedblright{}.) Say it once more.
\end{tcolorbox}
